\documentclass{article}
\usepackage{amsmath, amssymb, amsthm}

\setlength{\parindent}{0pt}

\newtheorem{claim}{Claim}

\begin{document}

Suppose there are four kinds of doughnuts: plain, chocolate, glazed, and butterscotch. 

\begin{claim}
    The generating function for the number of ways to select the flavours of $n$ doughnuts is
    \[
    \left(\sum_{k=0}^{\infty}x^{2k+1}\right)^4
    \]
    where each flavour occurs an odd number of times.
\end{claim}

\begin{proof}
    For a single flavour, the possible numbers of doughnuts are
    \[
    1, 3, 5, \dots
    \]
    so the corresponding generating function is
    \[
    x + x^3 + x^5 + \dots = \sum_{k=0}^{\infty}x^{2k+1}
    \]

    \medskip

    Since the counts for each flavour are chosen independently, the product rule for 
    generating functions gives
    \[
    \left(\sum_{k=0}^{\infty}x^{2k+1}\right)^4
    \]

    \medskip

    Hence the generating function is
    \[
    \left(\sum_{k=0}^{\infty}x^{2k+1}\right)^4
    \]
\end{proof}

\begin{claim}
    The generating function for the number of ways to select the flavours of $n$ doughnuts is 
    \[
    \left(\sum_{k=0}^{\infty}x^{3k}\right)^4
    \]
    where each flavour occurs a multiple of $3$ times.
\end{claim}

\begin{proof}
    For a single flavour, the possible number of doughnuts are
    \[
    0, 3, 6, \dots
    \]
    so the corresponding generating function is
    \[
    1 + x^3 + x^6 + \dots = \sum_{k=0}^{\infty}x^{3k}
    \]

    \medskip

    Since the counts for each flavour are chosen independently, the product rule for
    generating functions gives
    \[
    \left(\sum_{k=0}^{\infty}x^{3k}\right)^4
    \]

    \medskip

    Hence the generating function is
    \[
    \left(\sum_{k=0}^{\infty}x^{3k}\right)^4
    \]
\end{proof}

\begin{claim}
    The generating function for the number of ways to select the flavours of $n$ doughnuts is 
    \[
    \left(\sum_{k=0}^{\infty}x^k\right)^2\left(\sum_{k=0}^{0}x^k\right)
    \left(\sum_{k=0}^{1}x^k\right)
    \]
    where there are no chocolate doughnuts and at most one glazed doughnut.
\end{claim}

\begin{proof}
    For both plain and butterscotch, the possible number of doughnuts are
    \[
    0, 1, 2, \dots
    \]
    so the corresponding generating function is
    \[
    1 + x + x^2 + \dots = \sum_{k=0}^{\infty}x^k
    \]

    \medskip

    For chocolate, the only possible number of doughnuts is $0$, so the corresponding
    generating function is
    \[
    1 = \sum_{k=0}^{0}x^k
    \]
    For glazed, the possible number of doughnuts are $0$ or $1$, so the corresponding
    generating function is
    \[
    1 + x = \sum_{k=0}^{1}x^k
    \]

    \medskip

    Since the counts for each flavour are chosen independently, the product rule for
    generating functions gives
    \[
    \left(\sum_{k=0}^{\infty}x^k\right)^2\left(\sum_{k=0}^{0}x^k\right)
    \left(\sum_{k=0}^{1}x^k\right)
    \]

    \medskip

    Hence the generating function is
    \[
    \left(\sum_{k=0}^{\infty}x^k\right)^2\left(\sum_{k=0}^{0}x^k\right)
    \left(\sum_{k=0}^{1}x^k\right)    
    \]
\end{proof}

\begin{claim}
    The generating function for the number of ways to select the flavours of $n$ doughnuts is 
    \[
    \left(\sum_{k=0}^{\infty}x^k\right)^2(x + x^3 + x^{11})(x^2 + x^4 + x^5)
    \]
    where there are $1$, $3$, or $11$ chocolate and $2$, $4$, or $5$ glazed.
\end{claim}

\begin{proof}
    For both plain and butterscotch, the possible number of doughnuts are
    \[
    0, 1, 2, \dots
    \]
    so the corresponding generating function is
    \[
    1 + x + x^2 + \dots = \sum_{k=0}^{\infty}x^k
    \]

    \medskip

    For chocolate, the possible number of doughnuts are
    \[
    1, 3, 11
    \]
    so the corresponding generating function is
    \[
    x + x^3 + x^{11}
    \]
    Similarly, for glazed, the possible number of doughnuts are
    \[
    2, 4, 5
    \]
    so the corresponding generating function is
    \[
    x^2 + x^4 + x^5
    \]

    \medskip

    Since the counts for each flavour are chosen independently, the product rule for
    generating functions gives
    \[
    \left(\sum_{k=0}^{\infty}x^k\right)^2(x + x^3 + x^{11})(x^2 + x^4 + x^5)
    \]

    \medskip

    Hence the generating function is
    \[
    \left(\sum_{k=0}^{\infty}x^k\right)^2(x + x^3 + x^{11})(x^2 + x^4 + x^5)
    \]
\end{proof}

\begin{claim}
    The generating function for the number of ways to select the flavours of $n$ doughnuts is 
    \[
    \left(\sum_{k=10}^{\infty}x^k\right)^4
    \]
    where each flavour occurs at least $10$ times.
\end{claim}

\begin{proof}
    For a single flavour, the possible numbers of doughnuts are 
    \[
    10, 11, 12, \dots
    \]
    so the corresponding generating function is 
    \[
    x^{10} + x^{11} + x^{12} + \dots = \sum_{k=10}^{\infty}x^k
    \]

    \medskip
    
    Since the counts for each flavour are chosen independently, the product rule for
    generating functions gives
    \[
    \left(\sum_{k=10}^{\infty}x^k\right)^4
    \]

    \medskip

    Hence the generating function is
    \[
    \left(\sum_{k=10}^{\infty}x^k\right)^4
    \]
\end{proof}

\end{document}
